\section{Conclusion}
\label{sec:conclusion}

Declination occupies a distinctive position in the L-series: it is the metric with
the best geometric transparency, the largest discrimination between compact neutral
maps and deliberately cracked enacted maps (70--80\% reduction from enacted to
\textsc{bisect}), and the worst legal track record (zero court adoptions as of 2026).
These properties are related. The geometric sensitivity that makes Declination a
powerful discrimination tool also makes it vulnerable to the three methodological
critiques that have prevented judicial adoption.

The core empirical findings are clear: enacted maps in NC, WI, TX, and FL produce
Declination values of $+0.28$ to $+0.40$, while \textsc{bisect} algorithmic maps
produce $+0.07$ to $+0.09$. The gap of $0.21$--$0.31$ represents deliberate
packing of Democratic voters into fewer, higher-margin districts---a pattern that
the NC district-level decomposition (Table~\ref{tab:decl-nc-decomp}) shows is
driven primarily by elevated Democratic-won district centroids ($\bar{v}_D^+ \approx
0.71$ enacted vs.\ $0.64$ bisect) rather than by cracking of Republican-won margins
(similar $\bar{v}_R^-$ across both plans).

For expert witnesses, the recommended approach is to use Declination's geometric
intuition without leading on the scalar metric. The sorted vote-share plot that
underlies Declination is an effective courtroom exhibit in its own right: the visual
contrast between enacted maps (Republican-won districts clustered near 50\%, Democratic-won
districts far from 50\%) and \textsc{bisect} maps (both parties' won districts similarly
close to 50\%) communicates the packing-cracking structure directly. This visual
exhibit can be presented alongside EG, MM, and Partisan Bias as convergent evidence
without requiring the court to adopt Declination as a formal legal standard.

The question of whether courts will eventually adopt Declination as a metric is not
one this paper attempts to answer. The current legal status is clear: no court has
done so as of 2026. Future developments in state court redistricting doctrine may
change this, particularly if the methodological critiques are addressed through
calibration to a reference distribution of neutral maps---which the \textsc{bisect}
distribution can provide.

\bigskip

\noindent\textit{The author thanks the redistricting research team for data
access and algorithmic development. All errors are the author's own.}
